% =========================
% L24.tex
% =========================

\chapter{L24 -- Sintesi via Lyapunov, cicli limite, teoremi di instabilit\`a e metodo indiretto}
\label{chap:L24}

\section{Richiamo: progetto di un controllo tramite Lyapunov}

Riprendiamo l'idea vista nella parte finale della lezione precedente: la teoria di Lyapunov non serve solo per \emph{analizzare} la stabilit\`a di un sistema, ma pu\`o essere usata anche per \emph{progettare} un controllo stabilizzante.

Consideriamo il sistema non lineare controllato
\begin{equation}
\label{eq:L24-sistema-controllato}
\begin{cases}
\dot x_1 = (x_2+1)^3,\\[4pt]
\dot x_2 = (x_1-1)^3 + u + 2.
\end{cases}
\end{equation}

Vogliamo costruire un controllo \(u(t)\) tale che il punto
\[
(\bar x_1,\bar x_2)=(1,-1)
\]
sia un equilibrio stabile.

\subsection{Scelta dell'equilibrio e traslazione}

Per prima cosa osserviamo che, se fissiamo
\[
\bar u=-2,
\]
allora il punto \((1,-1)\) \`e effettivamente un equilibrio del sistema, perch\'e
\[
(x_2+1)^3 = (-1+1)^3 = 0,
\]
e inoltre
\[
(x_1-1)^3 + \bar u + 2 = (1-1)^3 -2 + 2 = 0.
\]

Per lavorare pi\`u comodamente trasliamo il sistema nell'origine:
\[
z_1 = x_1-1,
\qquad
z_2 = x_2+1,
\qquad
v = u-\bar u = u+2.
\]
Equivalentemente,
\[
u = v-2.
\]

In queste nuove coordinate il sistema diventa
\begin{equation}
\label{eq:L24-sistema-traslato}
\begin{cases}
\dot z_1 = z_2^3,\\[4pt]
\dot z_2 = z_1^3 + v.
\end{cases}
\end{equation}
Ora l'equilibrio desiderato coincide con l'origine \(z=0\).

\subsection{Scelta della funzione di Lyapunov}

Proviamo con la candidata naturale
\[
V(z)=\frac{1}{2}(z_1^2+z_2^2),
\]
che \`e positiva definita.

La derivata lungo le traiettorie di \eqref{eq:L24-sistema-traslato} \`e
\[
L_fV(z)=\frac{\partial V}{\partial z}\dot z
= z_1\dot z_1 + z_2\dot z_2.
\]
Sostituendo:
\begin{align*}
L_fV(z)
&= z_1 z_2^3 + z_2(z_1^3+v)\\
&= z_1 z_2^3 + z_1^3 z_2 + z_2 v.
\end{align*}

I primi due termini hanno segno indefinito. L'idea allora \`e usare il controllo per cancellarli e aggiungere un termine dissipativo.

Scegliamo
\[
v = - z_1 z_2^2 - z_1^3 - z_2.
\]
Allora
\begin{align*}
L_fV(z)
&= z_1 z_2^3 + z_1^3 z_2 + z_2\bigl(- z_1 z_2^2 - z_1^3 - z_2\bigr)\\
&= z_1 z_2^3 + z_1^3 z_2 - z_1 z_2^3 - z_1^3 z_2 - z_2^2\\
&= -z_2^2 \le 0.
\end{align*}
Quindi il sistema in anello chiuso \`e stabile.

\subsection{Legge di controllo in coordinate originali}

Ricordando che \(u=v-2\), otteniamo la legge di controllo
\[
\boxed{
u(x)=-(x_1-1)(x_2+1)^2-(x_1-1)^3-(x_2+1)-2.
}
\]

\medskip

\noindent
\textbf{Osservazione importante.}
Questo esempio mostra bene l'idea fondamentale:
\begin{center}
\emph{si pu\`o progettare il controllo imponendo direttamente il segno della derivata di una funzione di Lyapunov.}
\end{center}

\section{Controllo affine in \(x\)}

Una legge di controllo si dice \textbf{affine in \(x\)} se ha la forma
\[
u = Kx + h,
\]
oppure, in componenti,
\[
u = k_1 x_1 + k_2 x_2 + \cdots + k_n x_n + h.
\]

Si tratta quindi di un controllo lineare nelle variabili di stato, pi\`u una costante.

\medskip

\noindent
\textbf{Perch\'e compare la costante?}  
Perch\'e in generale l'equilibrio da stabilizzare non coincide con l'origine. Se vogliamo stabilizzare un punto \(\bar x\), allora il controllo deve soddisfare
\[
u(\bar x)=\bar u,
\]
cio\`e deve assumere il valore corretto nel punto di equilibrio. In coordinate traslate, questo equivale a chiedere che il controllo si annulli nell'origine:
\[
v(0)=0.
\]

\medskip

\noindent
In altre parole:
\begin{center}
\emph{un controllo stabilizzante non deve spostare il punto di equilibrio che vogliamo mantenere.}
\end{center}

\section{Sistema non lineare controllato vs linearizzato controllato}

Questa parte \`e concettualmente molto importante.

Sia dato il sistema non lineare controllato
\[
\dot x = f(x,u),
\]
e sia \((\bar x,\bar u)\) una coppia di equilibrio, cio\`e
\[
f(\bar x,\bar u)=0.
\]

Trasliamo le variabili:
\[
z = x-\bar x,
\qquad
v = u-\bar u.
\]

Sviluppando \(f(x,u)\) al primo ordine attorno a \((\bar x,\bar u)\), otteniamo
\[
f(x,u)\approx f(\bar x,\bar u)
+ \left.\frac{\partial f}{\partial x}\right|_{(\bar x,\bar u)} (x-\bar x)
+ \left.\frac{\partial f}{\partial u}\right|_{(\bar x,\bar u)} (u-\bar u).
\]
Poich\'e \(f(\bar x,\bar u)=0\), segue
\[
\dot z \approx Az + Bv,
\]
dove
\[
A=\left.\frac{\partial f}{\partial x}\right|_{(\bar x,\bar u)},
\qquad
B=\left.\frac{\partial f}{\partial u}\right|_{(\bar x,\bar u)}.
\]

Supponiamo ora di applicare una retroazione statica affine
\[
u = Kx + h.
\]
Affinch\'e \((\bar x,\bar u)\) resti equilibrio deve valere
\[
\bar u = K\bar x + h.
\]
Dunque, nelle variabili traslate,
\[
v = u-\bar u = K(x-\bar x)=Kz.
\]
Sostituendo nel linearizzato:
\[
\dot z \approx Az + BKz = (A+BK)z.
\]

\subsection{Conclusione concettuale}

Il \textbf{linearizzato del sistema non lineare gi\`a controllato} coincide, al primo ordine, con il \textbf{sistema linearizzato non controllato a cui si applica poi la stessa retroazione lineare}.

\medskip

\noindent
In simboli:
\[
\text{linearizza}\bigl(f(x,Kx+h)\bigr)
\quad \Longleftrightarrow \quad
\dot z = (A+BK)z.
\]

Questa osservazione spiega perch\'e, localmente, il progetto del controllo sul modello linearizzato \`e significativo anche per il sistema non lineare.

\medskip

\noindent
\textbf{Messaggio da ricordare.}
Se il linearizzato in anello chiuso
\[
\dot z=(A+BK)z
\]
\`e asintoticamente stabile, allora il sistema non lineare controllato \`e localmente asintoticamente stabile attorno all'equilibrio, per il metodo indiretto di Lyapunov.

\section{Ciclo limite}

Nei sistemi non lineari possono comparire fenomeni che nei sistemi lineari non si osservano in questo modo. Il pi\`u importante \`e il \textbf{ciclo limite}.

\subsection{Definizione qualitativa}

Un ciclo limite \`e una traiettoria periodica isolata nello spazio delle fasi.  
Questo significa che:

\begin{itemize}
    \item la traiettoria \`e chiusa;
    \item il sistema vi ritorna periodicamente;
    \item le traiettorie vicine possono essere attratte oppure respinte da tale curva.
\end{itemize}

\medskip

\noindent
\textbf{Attenzione:} un ciclo limite \emph{non} \`e un equilibrio.  
Un equilibrio \`e un punto fermo; un ciclo limite \`e un moto periodico.

\subsection{Idea per studiarne l'attrattivit\`a}

Se una curva chiusa \`e descritta implicitamente da
\[
C(x)=0,
\]
si pu\`o provare a costruire
\[
V(x)=\bigl(C(x)\bigr)^2.
\]
Questa funzione:

\begin{itemize}
    \item vale zero sulla curva;
    \item \`e positiva fuori dalla curva;
    \item misura, in un certo senso, la distanza energetica dalla curva.
\end{itemize}

La derivata lungo le traiettorie \`e
\[
L_fV(x)=2\,C(x)\,\frac{\partial C}{\partial x}f(x).
\]
Il segno di \(L_fV\) ci dice se le traiettorie si avvicinano alla curva oppure se ne allontanano.

\begin{itemize}
    \item Se \(L_fV<0\) vicino alla curva, il ciclo \`e attrattivo.
    \item Se \(L_fV>0\) vicino alla curva, il ciclo \`e repulsivo.
\end{itemize}

\subsection{Esempio classico}

Consideriamo il sistema
\begin{equation}
\label{eq:L24-ciclo}
\begin{cases}
\dot x_1 = -x_2 + x_1(1-x_1^2-x_2^2),\\[4pt]
\dot x_2 = x_1 + x_2(1-x_1^2-x_2^2).
\end{cases}
\end{equation}

\subsubsection{Studio dell'origine}

Con
\[
V(x)=\frac12(x_1^2+x_2^2),
\]
si ottiene
\begin{align*}
L_fV
&= x_1\dot x_1 + x_2\dot x_2\\
&= x_1\bigl(-x_2 + x_1(1-r^2)\bigr)
+ x_2\bigl(x_1 + x_2(1-r^2)\bigr),
\end{align*}
dove
\[
r^2=x_1^2+x_2^2.
\]
I termini misti si cancellano:
\[
L_fV = (x_1^2+x_2^2)(1-x_1^2-x_2^2)=r^2(1-r^2).
\]

Vicino all'origine vale \(r<1\), quindi
\[
L_fV>0
\qquad \text{per } x\neq 0 \text{ sufficientemente vicino all'origine}.
\]
Dunque l'origine \`e instabile.

\subsubsection{Studio della circonferenza \(x_1^2+x_2^2=1\)}

Poniamo
\[
C(x)=1-x_1^2-x_2^2 = 1-r^2,
\qquad
V_C(x)=\bigl(C(x)\bigr)^2.
\]
Allora
\[
\dot C = -2x_1\dot x_1 -2x_2\dot x_2.
\]
Sostituendo \eqref{eq:L24-ciclo}:
\begin{align*}
\dot C
&=
-2x_1\bigl(-x_2 + x_1(1-r^2)\bigr)
-2x_2\bigl(x_1 + x_2(1-r^2)\bigr)\\
&=
-2(x_1^2+x_2^2)(1-r^2)\\
&=
-2r^2 C.
\end{align*}
Di conseguenza
\[
L_fV_C = \frac{d}{dt}(C^2)=2C\dot C = -4r^2 C^2 \le 0.
\]
Fuori dall'origine e fuori dalla circonferenza unitaria, la derivata \`e strettamente negativa. Dunque la curva
\[
x_1^2+x_2^2=1
\]
\`e un ciclo limite attrattivo.

\begin{figure}[h!]
\centering
\begin{tikzpicture}[scale=1.2]
    \draw[->] (-2,0) -- (2,0) node[right] {$x_1$};
    \draw[->] (0,-2) -- (0,2) node[above] {$x_2$};

    \draw[thick,red] (0,0) circle (1);
    \node[red] at (1.35,1.2) {\small ciclo limite};

    \draw[orange,dashed,thick] (0,0) circle (0.55);
    \draw[orange,dashed,thick] (0,0) circle (1.45);

    \draw[-{Latex[length=3mm]},orange,thick] (0.55,0) -- (0.8,0);
    \draw[-{Latex[length=3mm]},orange,thick] (-0.55,0) -- (-0.8,0);
    \draw[-{Latex[length=3mm]},orange,thick] (0,0.55) -- (0,0.8);
    \draw[-{Latex[length=3mm]},orange,thick] (0,-0.55) -- (0,-0.8);

    \draw[-{Latex[length=3mm]},orange,thick] (1.45,0) -- (1.15,0);
    \draw[-{Latex[length=3mm]},orange,thick] (-1.45,0) -- (-1.15,0);
    \draw[-{Latex[length=3mm]},orange,thick] (0,1.45) -- (0,1.15);
    \draw[-{Latex[length=3mm]},orange,thick] (0,-1.45) -- (0,-1.15);

    \fill (0,0) circle (1.4pt);
    \node at (-0.45,-0.2) {\small origine instabile};
\end{tikzpicture}
\caption{Schema qualitativo: l'origine \`e instabile e la circonferenza unitaria \`e un ciclo limite attrattivo.}
\end{figure}

\section{Teoremi inversi di Lyapunov: idea generale}

I teoremi diretti di Lyapunov dicono:
\begin{center}
\emph{se trovi una funzione \(V\) con certe propriet\`a, allora puoi concludere qualcosa sulla stabilit\`a.}
\end{center}

I \textbf{teoremi inversi} affermano, in sostanza, il viceversa:

\begin{itemize}
    \item se un equilibrio \`e stabile, allora esiste una funzione di Lyapunov positiva definita con derivata non positiva;
    \item se un equilibrio \`e asintoticamente stabile, allora esiste una funzione di Lyapunov positiva definita con derivata negativa definita.
\end{itemize}

\medskip

\noindent
Questi risultati sono molto profondi, ma in pratica sono poco costruttivi:
\begin{center}
\emph{garantiscono che la funzione \(V\) esiste, ma non ti dicono come trovarla.}
\end{center}

Per questo, negli esercizi, il problema resta sempre lo stesso: bisogna saper scegliere una buona candidata.

\section{Teorema di instabilit\`a di Chetaev}
\label{sec:L24-chetaev}

Quando vogliamo dimostrare che un equilibrio \`e instabile, esiste un criterio molto utile: il teorema di Chetaev.

\subsection{Enunciato}

Consideriamo il sistema
\[
\dot x=f(x),
\]
con equilibrio nell'origine.  
Sia \(W\) un intorno dell'origine e sia \(V\in C^1(W)\). Supponiamo che esista un aperto \(U\subset W\) tale che:

\begin{enumerate}
    \item l'origine appartiene al bordo di \(U\);
    \item \(V(x)>0\) per ogni \(x\in U\);
    \item \(V(x)=0\) per ogni \(x\in \partial U \cap W\);
    \item \(L_fV(x)>0\) per ogni \(x\in U\).
\end{enumerate}

Allora l'origine \`e \textbf{instabile}.

\medskip

\noindent
L'idea geometrica \`e molto semplice: se in una regione \(U\) arbitrariamente vicina all'origine la funzione \(V\) cresce lungo le traiettorie, allora una traiettoria che parte in \(U\) non pu\`o rimanere confinata vicino all'origine.

\begin{figure}[h!]
\centering
\begin{tikzpicture}[scale=1.1]
    \draw[->] (-2.2,0) -- (2.2,0) node[right] {$x_1$};
    \draw[->] (0,-2.2) -- (0,2.2) node[above] {$x_2$};

    \draw[gray,thick] (0,0) circle (1.4);
    \node[gray] at (-1.6,1.6) {\small \(W\)};

    \fill[pattern=north east lines,pattern color=blue!60] (0,0) -- (1.4,0) arc (0:60:1.4) -- cycle;
    \draw[red,thick] (0,0) -- (1.4,0);
    \draw[red,thick] (0,0) -- ({1.4*cos(60)},{1.4*sin(60)});
    \draw[red,thick] (1.4,0) arc (0:60:1.4);
    \node at (0.8,0.45) {\small \(U\)};

    \fill (0,0) circle (1.5pt);
    \node at (-0.2,-0.2) {\small \(0\)};
\end{tikzpicture}
\caption{Schema qualitativo del teorema di Chetaev: l'origine sta sul bordo della regione \(U\).}
\end{figure}

\subsection{Esempio}

Consideriamo il sistema
\[
\begin{cases}
\dot x_1 = x_2^3,\\[4pt]
\dot x_2 = x_1^3.
\end{cases}
\]
L'origine \`e un equilibrio.

Scegliamo
\[
V(x)=x_1x_2.
\]
Allora
\[
V(0)=0,
\]
e nella regione del primo quadrante
\[
U=\{(x_1,x_2):x_1>0,\ x_2>0,\ x_1^2+x_2^2<\rho^2\}
\]
si ha
\[
V(x)>0.
\]
La derivata lungo le traiettorie \`e
\[
L_fV = x_2 \dot x_1 + x_1 \dot x_2 = x_2 x_2^3 + x_1 x_1^3 = x_1^4 + x_2^4 > 0
\]
per ogni \(x\in U\setminus\{0\}\). Quindi, per Chetaev, l'origine \`e instabile.

\section{Metodo indiretto di Lyapunov per sistemi lineari e non lineari}

Questa parte \`e centrale perch\'e spiega come il linearizzato permette di dedurre la stabilit\`a locale del sistema non lineare.

\subsection{Caso lineare continuo: equazione di Lyapunov}

Consideriamo il sistema lineare tempo continuo
\[
\dot x = Ax.
\]
Cerchiamo una candidata quadratica:
\[
V(x)=x^\top P x,
\]
con \(P=P^\top\).

La derivata lungo le traiettorie \`e
\begin{align*}
\dot V(x)
&= \dot x^\top P x + x^\top P \dot x\\
&= x^\top A^\top P x + x^\top P A x\\
&= x^\top (A^\top P + PA)x.
\end{align*}

Se imponiamo
\[
A^\top P + PA = -Q,
\]
con \(Q=Q^\top>0\), otteniamo
\[
\dot V(x) = -x^\top Q x <0 \qquad \forall x\neq 0.
\]

Questa \`e la \textbf{equazione algebrica di Lyapunov}.

\subsection{Teorema classico}

Dato \(Q=Q^\top>0\), esiste un'unica soluzione simmetrica \(P=P^\top>0\) dell'equazione
\[
A^\top P + PA = -Q
\]
se e solo se tutti gli autovalori di \(A\) hanno parte reale strettamente negativa.

Quindi:
\[
A \text{ Hurwitz}
\quad \Longleftrightarrow \quad
\exists\, P>0 \text{ t.c. } A^\top P + PA = -Q.
\]

\medskip

\noindent
Una rappresentazione utile della soluzione \`e
\[
P=\int_0^{+\infty} e^{A^\top t} Q e^{At}\,dt.
\]

\subsection{Caso lineare discreto}

Per il sistema
\[
x(k+1)=Ax(k),
\]
poniamo ancora
\[
V(x)=x^\top P x.
\]
Allora
\begin{align*}
\Delta V(x)
&=V(x(k+1))-V(x(k))\\
&=x(k)^\top A^\top P A x(k) - x(k)^\top P x(k)\\
&=x(k)^\top (A^\top P A - P)x(k).
\end{align*}
Se imponiamo
\[
A^\top P A - P = -Q,
\qquad Q>0,
\]
si ottiene
\[
\Delta V(x) = -x^\top Q x <0.
\]

Anche qui vale il teorema analogo:
\[
\exists\, P>0 \text{ t.c. } A^\top P A - P = -Q
\quad \Longleftrightarrow \quad
|\lambda_i(A)|<1 \ \forall i.
\]

Una formula utile per \(P\) \`e
\[
P = \sum_{k=0}^{+\infty} (A^\top)^k Q A^k.
\]

\section{Metodo indiretto per sistemi non lineari}

Consideriamo ora un sistema non lineare
\[
\dot x = f(x),
\qquad
f(0)=0.
\]
Supponiamo che \(f\) sia regolare e sviluppabile attorno all'origine:
\[
f(x)=Ax+h(x),
\qquad
A=\left.\frac{\partial f}{\partial x}\right|_{x=0},
\]
dove il resto \(h(x)\) soddisfa
\[
\lim_{x\to 0}\frac{\|h(x)\|}{\|x\|}=0.
\]

\subsection{Caso stabile del linearizzato}

Se \(A\) \`e Hurwitz, scegliamo \(Q=I\) e risolviamo
\[
A^\top P + PA = -I,
\]
ottenendo \(P>0\). Consideriamo allora
\[
V(x)=x^\top P x.
\]
Lungo le traiettorie del sistema non lineare:
\begin{align*}
L_fV(x)
&= x^\top(PA+A^\top P)x + 2x^\top P h(x)\\
&= -\|x\|^2 + 2x^\top P h(x).
\end{align*}
Per \(x\) sufficientemente piccolo, il termine \(h(x)\) \`e di ordine superiore e soddisfa una stima del tipo
\[
\|h(x)\| \le \frac{1}{4\|P\|}\|x\|.
\]
Di conseguenza
\[
2|x^\top P h(x)| \le 2\|x\|\,\|P\|\,\|h(x)\|
\le \frac12 \|x\|^2,
\]
e quindi
\[
L_fV(x)\le -\frac12\|x\|^2<0
\]
per \(x\neq 0\) sufficientemente vicino all'origine.

Allora l'origine del sistema non lineare \`e \textbf{localmente asintoticamente stabile}.

\subsection{Caso instabile del linearizzato}

Se \(A\) possiede almeno un autovalore con parte reale positiva, allora l'origine del sistema non lineare \`e instabile.

L'idea della dimostrazione \`e separare, tramite cambio di base, la parte instabile e quella stabile del linearizzato:
\[
A \sim
\begin{pmatrix}
A_1 & 0\\
0 & A_2
\end{pmatrix},
\]
dove:
\begin{itemize}
    \item \(A_1\) contiene gli autovalori con parte reale positiva;
    \item \(A_2\) contiene quelli con parte reale negativa.
\end{itemize}

Si costruisce quindi una funzione di Chetaev usando due soluzioni di Lyapunov sui due sottoblocchi e si dimostra che esiste una regione in cui la funzione cresce lungo le traiettorie. Da qui segue l'instabilit\`a dell'origine.

\subsection{Caso inconcludente}

Se tutti gli autovalori hanno parte reale \(\le 0\), ma qualcuno ha parte reale nulla, il metodo indiretto \textbf{non conclude}.  
In questo caso servono strumenti ulteriori: Lyapunov diretto, LaSalle, studio qualitativo, centro-manifold, eccetera.

\section{Sistemi non stazionari: perch\'e il criterio sugli autovalori non vale pi\`u}
\label{sec:L24-nonstazionari}

Questa parte \`e soprattutto concettuale.

Consideriamo un sistema lineare tempo-variante
\[
\dot x(t)=A(t)x(t).
\]
A differenza del caso stazionario, qui gli autovalori istantanei di \(A(t)\) \emph{non bastano} per dedurre la stabilit\`a del sistema.

\subsection{Esempio 1: autovalori sempre negativi ma sistema instabile}

Consideriamo
\[
A(t)=
\begin{pmatrix}
-1 & e^{2t}\\
0 & -1
\end{pmatrix}.
\]
Gli autovalori di \(A(t)\) sono sempre
\[
\lambda_1(t)=\lambda_2)=\lambda_2)=\lambda_2(t)=-1,
\]
quindi sembrerebbe un sistema stabile. Ma il sistema \`e
\[
\begin{cases}
\dot x_1 = -x_1 + e^{2t}x_2,\\[4pt]
\dot x_2 = -x_2.
\end{cases}
\]
Dalla seconda equazione:
\[
x_2(t)=e^{-t}x_2(0).
\]
Sostituendo nella prima:
\[
\dot x_1 + x_1 = e^{2t}e^{-t}x_2(0)=e^t x_2(0).
\]
Moltiplicando per \(e^t\):
\[
\frac{d}{dt}\bigl(e^t x_1(t)\bigr)=e^{2t}x_2(0).
\]
Integrando:
\[
e^t x_1(t)=x_1(0)+\frac{x_2(0)}{2}(e^{2t}-1),
\]
quindi
\[
x_1(t)=e^{-t}x_1(0)+\frac{x_2(0)}{2}(e^t-e^{-t}).
\]
Se \(x_2(0)\neq 0\), allora \(x_1(t)\) diverge come \(e^t\). Quindi il sistema \`e instabile, pur avendo autovalori istantanei sempre uguali a \(-1\).

\subsection{Messaggio da ricordare}

Per i sistemi non stazionari:
\begin{center}
\emph{gli autovalori ``congelati'' di \(A(t)\) non determinano da soli la stabilit\`a del sistema.}
\end{center}

Ci\`o che conta \`e l'evoluzione complessiva nel tempo, non il solo spettro istantaneo.

\section{Schema finale della lezione}

La lezione raccoglie diversi messaggi importanti.

\begin{enumerate}
    \item \textbf{Lyapunov pu\`o essere usato per la sintesi del controllo}: scegli una funzione \(V\) e costruisci \(u\) in modo che \(L_fV\le 0\) o \(L_fV<0\).
    
    \item \textbf{Il linearizzato controllato} del sistema non lineare coincide, al primo ordine, con il sistema \((A+BK)\) ottenuto applicando la retroazione al linearizzato.

    \item \textbf{I cicli limite} sono traiettorie periodiche isolate tipiche dei sistemi non lineari: non sono equilibri, ma possono essere attrattivi o repulsivi.

    \item \textbf{Il teorema di Chetaev} \`e uno strumento utile per dimostrare instabilit\`a quando si riesce a costruire una funzione che cresce in una regione arbitrariamente vicina all'origine.

    \item \textbf{L'equazione di Lyapunov} fornisce una procedura concreta per costruire candidate quadratiche nei sistemi lineari stabili.

    \item \textbf{Il metodo indiretto di Lyapunov} collega il segno degli autovalori del linearizzato con la stabilit\`a locale del sistema non lineare.

    \item \textbf{Nei sistemi tempo-varianti}, invece, il criterio sugli autovalori istantanei non \`e pi\`u affidabile.
\end{enumerate}

\medskip

\noindent
\textbf{Frase da ricordare all'orale.}
\begin{center}
\emph{Nel caso non lineare, il linearizzato descrive il comportamento locale solo vicino all'equilibrio; se il linearizzato \`e asintoticamente stabile, allora anche il sistema non lineare lo \`e localmente. Ma se il linearizzato \`e inconcludente, servono Lyapunov diretto, LaSalle o strumenti qualitativi pi\`u fini.}
\end{center}